\documentclass[11pt]{article}
\usepackage[margin=1in]{geometry}
\usepackage{amsmath, amssymb}
\usepackage{graphicx}
\usepackage{booktabs}
\usepackage{hyperref}

\title{Distilling a Floor-Plan Vision-Language Model into a Tiny VLM}
\author{Project Notes}
\date{}

\begin{document}
\maketitle

\begin{abstract}
We describe a reproducible pipeline to distill a large vision-language teacher into a compact student model for evidence-grounded floor-plan question answering. The method ingests raster floor plans, prompts a vision-capable teacher to produce answers and evidence, applies a critic filter, and trains a student in three stages: answer distillation, evidence distillation, and preference optimization for abstention. Evaluation covers answer accuracy, evidence quality, abstention behavior, robustness, and efficiency.
\end{abstract}

\section{Introduction}
Given floor-plan imagery, we seek a small VLM that can answer questions, point to supporting evidence (door/room IDs or masks), and abstain when uncertain. We leverage a stronger teacher (e.g., GPT-4o) to generate weak labels and distill knowledge into a lighter student (CLIP vision + small text decoder).

\section{Data}
\textbf{Input}: raster floor plans (e.g., CubiCasa5k). Each sample includes image $\mathbf{x}$, metadata (paths, rotation), and optional geometric annotations (rooms, doors).\\
\textbf{Labels}: JSONL with fields $(q, a, \mathcal{E}, u)$ where $q$ is question, $a$ is answer token/string, $\mathcal{E}$ contains evidence pointers (door/room IDs, optional mask), $u$ is uncertainty. We collect $k=3$ diverse questions per image from a large bank.

\section{Teacher Labeling}
\textbf{Prompt}: cautious reviewer; answer only if visible, otherwise \texttt{unknown}; return answer + evidence IDs + rationale.\\
\textbf{Generation}: For each image, sample $k$ questions, send low-res JPEG plus question to a vision-capable teacher (default GPT-4o-mini).\\
\textbf{Filtering}: Critic discards entries with missing evidence (when required), contradictory tool traces, empty/none answers, or ``unseen'' rationales.\\
\textbf{Traceability}: Each record stores \texttt{image\_id}, \texttt{image\_path}, \texttt{question\_key}, and \texttt{qid}.

\section{Student Architecture}
We fuse a pretrained vision tower with a small text decoder.\\
\textbf{Vision}: CLIP ViT (e.g., \texttt{openai/clip-vit-base-patch32}), often frozen.\\
\textbf{Text}: Small LM (e.g., \texttt{gpt2}, or Phi-2/TinyLlama variants), optionally LoRA-tuned.\\
\textbf{Connector}: Concatenate pooled vision embedding $\mathbf{v} \in \mathbb{R}^{d_v}$ and pooled text embedding $\mathbf{t} \in \mathbb{R}^{d_t}$, then
\[
\mathbf{h} = \phi(W[\mathbf{v};\mathbf{t}] + b), \quad \phi=\text{GELU},
\]
followed by heads for answer logits, evidence pointer logits, and an abstain logit. Evidence head outputs a fixed-dim pointer vector.

\section{Training Objectives}
Let $z_s$ be student answer logits, $z_t$ teacher logits, $y$ target token, $T$ temperature. Let $p_s, p_t$ be student/teacher evidence pointers, and $\beta$ the DPO temperature.
\subsection{Stage A: Answer KD}
\[
L_A = \lambda_{\text{ce}}\cdot \text{CE}(z_s, y) + \lambda_{\text{kd}}\cdot T^2 \cdot \text{KL}\big(\text{softmax}(z_t/T) \,\|\, \text{softmax}(z_s/T)\big).
\]
\subsection{Stage B: Evidence KD}
\[
L_B = L_A + \lambda_{\text{ev}}\cdot \|p_s - p_t\|_2^2.
\]
\subsection{Stage C: Abstention via DPO}
With preference log-probs $\pi$ (student) and $\rho$ (reference), positive = abstain-when-uncertain, negative = answer-when-uncertain,
\[
L_C = -\log \sigma\big(\beta\big[(\pi_{\text{pos}} - \pi_{\text{neg}}) - (\rho_{\text{pos}} - \rho_{\text{neg}})\big]\big).
\]

\section{Training Regimen}
\textbf{Stage A} (Answer KD): train on filtered labels for answer logits + KD.\\
\textbf{Stage B} (Evidence KD): initialize from Stage A; add evidence pointer loss.\\
\textbf{Stage C} (DPO): initialize from Stage B; optimize abstention behavior.\\
\textbf{Hyperparameters}: image size 224; max length 128; freeze vision; small LR for text ($\sim$2e-5) and projector ($\sim$1e-4); cosine schedule; gradient accumulation as needed.

\section{Evaluation}
\textbf{QA}: exact/lenient accuracy.\\
\textbf{Evidence}: F1/IoU on pointers/masks.\\
\textbf{Abstention}: precision/recall/AUPRC on \texttt{unknown} decisions.\\
\textbf{Graph}: adjacency edit distance (if room graphs available).\\
\textbf{Efficiency}: params, VRAM (bs=1), latency.\\
\textbf{Robustness}: rotation, scale, JPEG noise, crops.\\
\textbf{Cross-dataset}: train on CAD, test on raster, and vice versa.

\section{Design Choices}
\begin{itemize}
    \item Pretrained vision/text for strong priors and rapid convergence.
    \item Three-stage KD for stable learning and explicit evidence supervision.
    \item Multiple questions per image to maximize label efficiency.
    \item Traceability fields to audit per-image/per-question outputs.
    \item Abstention-first safety: prefer \texttt{unknown} when evidence is insufficient.
\end{itemize}

\section{Future Work}
\begin{itemize}
    \item Swap text to Phi-2/TinyLlama with QLoRA; unfreeze top vision blocks or add adapters.
    \item Integrate CAD tokens and geometry tools directly into the model input or retrieval.
    \item Heatmap-based evidence supervision instead of pointer-only.
    \item Richer preference pairs for DPO from tool traces and rationales.
\end{itemize}

\end{document}
